%------------------------------------------------------------------------------%
\section{Função do segundo grau ou função quadrática}

Uma função é dita
\emph{quadrática}\index{Função!quadrática}\index{Função!segundo grau}
se pode ser escrita na forma de um polinômio do 2\textsuperscript{o} grau
\[
  f(x) = ax^2 + bx + c,
\]
em que $a, b, c \in \R$ e $a\neq 0$.
O gráfico de uma função quadrática é uma parábola e é possível
esboçá-lo a partir de sua(s) raiz(es), caso exista(m),
ou das coordenadas do vértice.

%------------------------------------------------------------------------------%
\subsection{Encontrando zeros e as coordenadas do vértice da parábola}

Para que possamos deduzir os zeros de uma função quadrática,
vamos completar quadrados.
\begin{align*}
  ax^2+bx+c & = 0                                   & & \text{colocando $a$ em evidência} \\[2mm]
  a \left(x^2+\frac{bx}{a}+\frac{c}{a}\right) & = 0 & & \text{sendo } a \neq 0 \\[2mm]
  x^2+\frac{bx}{a}+\frac{c}{a} & = 0                & & \text{completando quadrado} \\[2mm]
  \left(x+\frac{b}{2a}\right)^2-\frac{b^2}{4a^2}+\frac{c}{a} & = 0 \\
  \left(x+\frac{b}{2a}\right)^2 & = \frac{b^2-4ac}{4a^2}
\end{align*}
Definindo
\[
  \Delta = b^2-4ac
\]
temos que a equação $ax^2+bx+c=0$ possui raiz real
se e somente se $\Delta \geq 0$ e nesse caso
\[
  x+\frac{b}{2a} = \pm\frac{\sqrt{\Delta}}{2a}.
\]

%------------------------------------------------------------------------------%
\begin{observation}
Resumindo.

Se $\Delta < 0$, a equação não possui raízes reais (as raízes são
complexas).

Se $\Delta > 0$, a equação possui duas raízes reais distintas:
\[
  x_1 = \frac{-b-\sqrt{\Delta}}{2a} \qquad\text{ e }\qquad
  x_2 = \frac{-b+\sqrt{\Delta}}{2a}.
\]

Se $\Delta = 0$, a equação possui uma única raiz real de
multiplicidade 2 (raiz dupla)
\[
  x_1 = \frac{-b}{2a}=x_2.
\]
\end{observation}
%------------------------------------------------------------------------------%

Em ambos os casos, $f(x)=ax^2+bx+c$ pode ser reescrita como
\[
  f(x) = a\left( x+\frac{b}{2a} \right)^2 - \frac{\Delta}{4a}.
\]
Logo, o gráfico de $f$ pode ser pensado como uma dupla
translação de
\[
  g(x)=ax^2.
\]
Transladamos
\[
  \frac{-\Delta}{4a}
\]
no eixo $y$ e
\[
  \frac{-b}{2a}
\]
no eixo $x$. Como $g$ possui vértice $(0, 0)$, $f$ possui vértice
\[
  (x_v, y_v) = \left(\frac{-b}{2a}, \frac{-\Delta}{4a}\right)
\]
e seu gráfico é simétrico em relação à reta $x=x_v$.

Se $a > 0$ então a parábola tem concavidade voltada para cima,
além disso $f$ tem um único ponto de mínimo dado por
\[
  x_v=\frac{-b}{2a}
\]
e o valor mínimo é dado por
\[
  y_v= \frac{-\Delta}{4a}.
\]
Se $a < 0$ então a parábola tem concavidade voltada para baixo,
além disso $f$ tem um único ponto de máximo dado
\[
  x_v=\frac{-b}{2a}
\]
e o valor máximo é dado por
\[
  y_v= \frac{-\Delta}{4a}.
\]
A Figura a seguir ilustra o que foi discutido nessa seção.

%------------------------------------------------------------------------------%
\begin{center}
  \begin{tabular}{|c|c|c|}
  \hline
   & $a>0$ & $a<0$ \\ \hline
  $\Delta >0$ & \includegraphics[scale=0.3]{sinq1} & \includegraphics[scale=0.3]{sinq2}  \\
  \hline
  $\Delta =0$ & \includegraphics[scale=0.3]{sinq3} & \includegraphics[scale=0.3]{sinq4}  \\
  \hline
  $\Delta <0$ & \includegraphics[scale=0.3]{sinq5} & \includegraphics[scale=0.3]{sinq6}  \\
  \hline
  \end{tabular}
\end{center}
%------------------------------------------------------------------------------%

%------------------------------------------------------------------------------%
\begin{observation}
  Se tivermos as coordenadas dos vértices, os zeros e o interceptos com o eixo
  $y$ podemo fazer o esboço do gráfico de uma equação do segundo grau com muita
  facilidade.
\end{observation}
%------------------------------------------------------------------------------%

%------------------------------------------------------------------------------%
\begin{example}
  Considere a seguinte função
  \[
    f(x)=x^2-5x+6.
  \]
  Encontre os zeros, as coordenadas do vértice, o intercepto com o eixo $y$
  e use esses dados para fazer o esboço do gráfico de $f$.

  \solution
  \emph{Zeros}

  Para encontrarmos os zeros precisamos resolver a equação
  \begin{align*}
    f(x) & =0 \\
    x^2-5x+6 & = 0,
  \end{align*}
  cujo
  \[
    \Delta = (-5)^2 - 4 \cdot 1 \cdot 6 = 1
  \]
  e
  \[
    x = \frac{-(-5) \pm \sqrt{1}}{2 \cdot 1}
  \]
  obtendo $x=2$ ou $x=3$, portanto temos os pontos $P_1=(2, 0)$ e
  $P_2=(3, 0)$, esses pontos são os interceptos com o eixo $x$.

  \emph{Vértices}

  \[
    x_v
    = \frac{-b}{2a}
    = \frac{-(-5)}{2 \cdot 1}
    = \frac{5}{2}
  \]
  e
  \[
    y_v
    = \frac{-\Delta}{4a}
    = \frac{-1}{4 \cdot 1}
    = -\frac{1}{4},
  \]
  logo
  \[
    V=\left(\frac{5}{2}, -\frac{1}{4}\right).
  \]

  \emph{Intercepto com o eixo $y$}

  Para encontrarmos o intercepto com o eixo $y$, temos que encontrar
  \[
    f(0) = 0^2 - 5 \cdots 0 + 6 = 6,
  \]
  logo temos o ponto $P_4=(0, 6)$.

  \emph{Gráfico}

  O gráfico dessa função está esboçado na imagem a seguir.
  \IncludeGraphicCentredEx{parab1}[0.5]
\end{example}
%------------------------------------------------------------------------------%

%------------------------------------------------------------------------------%
\begin{example}
  Faça o esboço do gráfico da parábola
  \(
    y = x^2 - x + 2
  \).

  \solution
  \emph{Zeros}

  Para encontrarmos os zeros precisamos resolver a equação
  \begin{align*}
    f(x)    & = 0 \\
    x^2-x+2 & = 0,
  \end{align*}
  cujo
  \[
    \Delta
    = (-1)^2-4 \cdot 1 \cdot 2
    = -7
    < 0,
  \]
  logo a função não possui zeros, com isso não intercepta o eixo $x$.

  \emph{Vértice}

  \[
    x_v
    = \frac{-b}{2a}
    = \frac{-(-1)}{2 \cdot 1}
    = \frac{1}{2}
  \]
  e
  \[
    y_v
    = \frac{-\Delta}{4a}
    = \frac{-(-7)}{4 \cdot 1}
    = \frac{7}{4},
  \]
  logo
  \[
    V = \left(\frac{1}{2}, \frac{7}{4}\right).
  \]

  \emph{Intercepto com o eixo $y$}

  Para encontrarmos o intercepto com o eixo $y$,
  temos que encontrar
  \[
    f(0) = 0^2 - 0 + 2 = 2,
  \]
  logo temos o ponto $P_1=(0, 2)$.

  Note que temos apenas 2 pontos para marcar no gráfico, para que possamos ter
  uma idea melhor da parábola vamos escolher outro ponto do domínio. Como
  já temos o ponto $P_1$, cuja coordenada $x$ está (na reta real) antes do
  $x_v$, vamos calcular $f(2)=4$ e daí temos o ponto $P_2=(2, 4)$.

  \emph{Gráfico}

  O gráfico dessa função está esboçado na imagem a seguir.
  \IncludeGraphicCentredEx{parab2}[0.5]
\end{example}
%------------------------------------------------------------------------------%

%------------------------------------------------------------------------------%
\begin{example}
  Faça o esboço do gráfico da parábola
  \(
    y=-x^2+4x
  \).

  \solution
  \emph{Zeros}

  Para encontrarmos os zeros precisamos resolver a equação
  \begin{align*}
    f(x) & =0 \\
    -x^2 + 4x = 0,
  \end{align*}
  como essa equação é incompleta ($c=0$), vamos resolver fatorando,
  ou seja,
  \[
    -x^2 + 4x = x(-x+4)
  \]
  portanto $x = 0$ ou $x = 4$.
  Assim, temos os ponto $P_1=(0, 0)$ e $P_2=(4, 0)$.

  \emph{Vértice}

  \[
    x_v
    = \frac{-b}{2a}
    = \frac{-4}{2(-1)}
    = 2
  \]
  e
  \[
    y_v
    =  \frac{-\Delta}{4a}
    = -\frac{16}{4(-1)}
    = 4,
  \]
  logo
  \[
    V=\left(2, 4\right).
  \]

  \emph{Intercepto com o eixo $y$}

  Para encontrarmos o intercepto com o eixo $y$, temos que encontrar
  \[
    f(0) = -0^2 + 4 \cdot 0 = 0,
  \]
  logo temos o ponto
  \[
    P_4 = P_1 = (0, 0).
  \]

  \emph{Gráfico}

  O gráfico dessa função está esboçado na imagem a seguir.
  \IncludeGraphicCentredEx{parab3}[0.5]
\end{example}
%------------------------------------------------------------------------------%

%------------------------------------------------------------------------------%